\begin{enumerate}
    \item \mbox{}
          \begin{minipage}[t][][t]{0.65\linewidth}
              On applique le théorème d'énergie cinétique (voir figure
              ci-contre)~:
          
              \[
                  \Delta E_c=\sum W_{AB}(\vec{F}_{\text{ext}})
              \]
              avec~:
              \[
                  \left\{
                  \begin{array}{l}
                      \Delta E_c=\frac{1}{2}m v_B^2-\frac{1}{2}m v_A^2=
                      0-\frac{1}{2}m v^2 \\[3mm]
                      \sum W_{AB}(\vec{F}_{\text{ext}})=W(\vec{P})=
                      mg(z_{A}-z_{B})=-mgh
                  \end{array}\right.
              \]
              donc~:
              \[
                  -\frac{1}{2}m v^2=-mgh
                  \quad\Rightarrow\
                  \mathbox{h=\frac{v^{2}}{2g}}
              \]
          
              Application numérique~:
              \[
                  h=\frac{(15,0)^{2}}{2\times 9,8}
                  \quad\Rightarrow\
                  h\approx \qty{11.5}{\m}
              \]
          
          \end{minipage}
          \hfill
          \begin{minipage}[t][][t]{0.34\linewidth}
              \begin{figure}[H]
                  \centering
                  \begin{tikzpicture}
                      \def\h{4}
                      \foreach \z/\n in {0/A,\h/B}{
                          \node[circle,ball color=black!40,inner sep=0pt,
                              minimum size=3mm,
                              label=left:$\n$] (\n) at (0,\z) {};
                          \node[circle,inner sep=0.6pt,
                              fill] (\n) at (\n.center) {};
                      }
                      \draw[<->,shorten <=1pt,shorten >=1pt]
                          (-1,0) -- node[right] {$h$} ++(0,\h);
                      \draw[->,thick,red] (A.center) -- ++(0,1)
                          node[right] {$\vec{v}$};
                      \draw[->,thick] (1,\h/3) --
                          node[right] {mouvement} ++(0,1);
                      \node[right=2mm] at (A.east) {$v_{A}=v=\qty{15.0}{\v}$};
                      \node[right=2mm] at (B.east) {$v_{B}=0$};
                  \end{tikzpicture}
              \end{figure}
          \end{minipage}
    \item \mbox{}
          \begin{minipage}[t][][t]{0.52\linewidth}
              On applique T.E.C entre $A$ et $C$ (voir figure ci-contre)~:
          
              \begin{align*}
                  &\frac{1}{2}m v_C^2-\frac{1}{2}m v_A^2=W_{AC}(\vec{P}) \\
                  &\Rightarrow \frac{1}{2}m v_C^2-\frac{1}{2}m v_A^2=-mgh^{\prime} \\
                  &\Rightarrow \frac{1}{2}m v_C^2=\frac{1}{2}m v_A^2-mgh^{\prime} \\
                  &\Rightarrow v_C^2=v_A^2-2gh^{\prime} \\
                  &\Rightarrow v_C=\sqrt{v_A^2-2gh^{\prime}}
              \end{align*}
          
              Application numérique~:
              \[
                  v_C=\sqrt{(15,0)^2-2\times 9,8\times 2}
                  \quad\Rightarrow\
                  v_{C}\approx \qty{13.6}{\v}
              \]
          \end{minipage}
          \hfill
          \begin{minipage}[t][][t]{0.46\linewidth}
              \begin{figure}[H]
                  \centering
                  \begin{tikzpicture}
                      \def\h{4}
                      \def\hp{1}
                      \foreach \z/\n in {0/A,\hp/C,\h/B}{
                          \node[circle,ball color=black!40,inner sep=0pt,
                              minimum size=3mm,
                              label=left:$\n$] (\n) at (0,\z) {};
                          \node[circle,inner sep=0.6pt,
                              fill] (\n) at (\n.center) {};
                      }
                      \draw[<->,shorten <=1pt,shorten >=1pt]
                          (-3,0) -- node[left] {$h=\qty{11.5}{\m}$} ++(0,\h);
                      \draw[<->,shorten <=1pt,shorten >=1pt]
                          (-1,0) -- node[left] {$h^{\prime}=\qty{2}{\m}$}
                          ++(0,\hp);
                      \draw[->,thick] (0.7,\h/2) --
                          node[right] {mouvement} ++(0,1);
                      \node[right=2mm] at (A.east) {$v_{A}=\qty{15.0}{\v}$};
                      \node[right=2mm] at (B.east) {$v_{B}=0$};
                      \node[right=2mm] at (C.east) {$v_{C}$};
                  \end{tikzpicture}
              \end{figure}
          \end{minipage}
    \item Soit $D$ le point d'arrivée au sol, donc~: $z_{A}=z_{D}$.
              
          On applique T.E.C entre $A$ et $D$~:
          \[
              \frac{1}{2}m v_D^2-\frac{1}{2}m v_A^2=mg(z_{A}-z_{D})
              \quad\Rightarrow\ \frac{1}{2}m v_D^2-\frac{1}{2}m v_A^2=0
              \quad\Rightarrow\ v_D=v_A=\qty{15.0}{\v}
          \]
\end{enumerate}

